\section{Review Questions}

\begin{enumerate}
\item Explain the transition from average rate of change to instantaneous rate of change. How does the geometry of secant and tangent lines represent this transition?

\item State the definition of the derivative at a point. Derive the derivative of $f(x)=x^2$ from the difference quotient and interpret the result.

\item Distinguish continuity from differentiability. Why does differentiability imply continuity, and what geometric behaviours can cause differentiability to fail?

\item State the elementary derivative rules used in the chapter. For each of the product, quotient, and chain rules, describe the structural pattern that the rule recognises.

\item Differentiate one low-degree polynomial, one elementary trigonometric function, and one simple composition. Keep the calculation short and explain the meaning of the resulting derivative in each case.

\item Explain how the sign of $f'$ determines monotonicity. What role does the mean value theorem play in passing from local derivative values to behaviour on an interval?

\item Define a critical point. Explain why $f'(c)=0$ identifies only a candidate for an extremum and how a sign chart distinguishes maxima, minima, and non-extremal critical points.

\item Explain what the second derivative reveals about concavity and changes of graph shape. What additional condition is required before a zero of $f''$ is called an inflection point?

\item Solve one elementary optimisation problem by following the full modelling workflow:
\begin{enumerate}
\item define the quantity to optimise;
\item reduce it to one variable;
\item determine the allowed domain;
\item find and compare all relevant candidates;
\item interpret the result in the original setting.
\end{enumerate}

\item Given qualitative information about $f'$ and $f''$, sketch a plausible graph of $f$. Mark intervals of increase and decrease, local extrema, concavity, and possible inflection points, and explain which conclusions are forced and which remain undetermined.
\end{enumerate}
